\section{Literature and Theory}
\label{sec:theory}

\todo{"recent studies" section for pumps and turbines, in particular, MOC is incredibly old and sudhof, erikson have modernized, state this.}

A note on the age of the methods used in this chapter. Much of the supersonic blading toolchain dates from the 1950s and 60s \cite{boxerApplicationSupersonicVortexflow1952, goldmanAnalyticalInvestigationSupersonic1968}, which can give the impression of outdated technology. This is not the case. The same methods were used in the development of the HM60 flight engine turbine \cite{erikssonDevelopmentSupersonicTurbine}, were given a modern doctoral treatment in 2020 \cite{sudhofDevelopmentTechniquesSupersonic2020}, and are still the basis of current small turbine design studies \cite{pengDesignAnalysis152023, abdalhamidDesignAnalysisSinglestage2024}. On the pump side, the Barske architecture is seeing renewed interest for electric rocket pumps \cite{danieliConceptualDesignElectric2023, knyazevaDesignFluidDynamics2021}.

\subsection{Pump Theory}

\subsubsection{Losses at Low Specific Speed}

\todo{I just realized I don't have the numbers to back me up that disc friction is better. I will waffle, then fill with literature and numbers later on.}

Centrifugal pumpss have multiple main sources of losses. The one that becomes significant at low specific speeds is the disc friction losses. When a rotor object rotates in a fluid, the shear stresses in the fluid cause a drag torque on the object. In most centrifugal pumps the impeller is disc shaped, and the loss is called disc friction loss. \todo{gulich pdf pg 144 says axial case cleareance only has a very small effect on disc friction. It also says that due to leakage flow, partload recirculation and exchange of momentum, tolerance of calculated disc friction loss can be at least +/-25\%, also, somehow head coefficient is related to disc friction? it says to maximize it for low specific speed}. With low specific speed pumps, disk frction is the main source of loss. The disc friction power of a rotating disc scales as

\begin{equation}
    P_{RR} = k_{RR}\, \rho\, \omega^3 r_2^5
    \label{eq:disc_friction_scaling}
\end{equation}

where $k_{RR}$ is a friction coefficient that falls slowly with Reynolds number \cite{gulichCentrifugalPumps2020}. The useful power scales with $\rho g Q H$, so as specific speed falls and the impeller diameter grows relative to the flow, the ratio of disc friction to useful power blows up. At $n_q=10$, $\psi_{\textrm{opt}=1}$ and $\mathrm{Re}=10^7$, the disc friction is approxiomately $30\%$ of the usefulf power of the impeller \cite{gulichpage103}. Gulich suggests that a high head coefficient \todo{high head coeffient not defined yet} reduces the diameter and consequently less disk frciont since $P_{RR}\sim D^5$. This is the reason why low specific speed centrifugal pumps physically cannot achieve the same high efficiencies as higher specific speed pumps. \todo{this paragraph is almost fully copied from gulich, need to rephrase.}

\todo{gulich page 429 says impeller losses are low at low specific speeds, referencing fig 3.33. Oop, fig3.33 shows the disc friction loss i think?}On the other hand, at low specific speeds, impeller losses are low as shown by \cref{fig3.33gulich}. \todo{explain this graph} Conversely, collector losses become dominant. \todo{need to explain what a collector is} Gulich explains \cite{gulich7.12} that for $n_q<12$ even higher efficiencies are expected than with volutes. This simplifies the design and manufacturing of the pump, which is a significant advantage at small scales where tight tolerances are hard to hold. \todo{need to explain what a volute is} This is important as the high fluid velocity at the peripheri makes the surface finish of the annular casing of primary importance for pump efficiency \cite{lockdesignfeatures}.

\subsubsection{The Barske Partial-Emission Pump}

\begin{figure}[H]
    \centering
    \begin{subfigure}[b]{0.45\textwidth}
        \centering
        
\includegraphics[width=\textwidth]{placeholder.png}
        \caption{Conventional shrouded full emission impeller} % PLACEHOLDER IMAGE
        \label{fig:impeller-conventional}
    \end{subfigure}
    \hfill
    \begin{subfigure}[b]{0.45\textwidth}
        \centering
        
\includegraphics[width=\textwidth]{placeholder.png}
        \caption{Barske open partial emission impeller} % PLACEHOLDER IMAGE
        \label{fig:impeller-barske}
    \end{subfigure}
    \caption{Comparison of the conventional and Barske pump architectures. The Barske impeller is open, has straight radial blades, and discharges into an annular casing with a single tangential diffuser.}
    \label{fig:impeller-comparison}
\end{figure}

\begin{figure}[H]
    \centering
    
\includegraphics[width=0.55\textwidth]{placeholder.png}
    \caption{Breakdown of pump losses against specific speed, after Gülich. At low specific speed the disc friction dominates while the impeller (blade) losses are small.} % PLACEHOLDER IMAGE: Gulich Fig 3.33 reproduction
    \label{fig3.33gulich}
\end{figure}

The Barske impeller pump design which was developed for the need for manufacturing simplicity and low specific speed applications became one of the highest efficiency designs at this specific speed. \todo{i might want a section on explaining basic pump and turbine geometry and terminology before this.} The barske impeller operates differently to a normal full emission pump. In a conventional centrifugal pump, impeller blade geometry aims to perfectly guide fluid flow along the blade surfaces in a smooth spiral streamline out of the impeller into the volute \cite{something}. The barske impeller is instead designed to induced almost a pure forced vortex flow where $u=\omega r$ and the fluid is essentially rotating as a solid body with the impeller \cite{barskeunconventionalpaper}. This also explains the unimportance of surface finish on the blades, as low relative velocity between the fluid and the impeller means lower friction losses. 

\todo{I know that conventional pumps are impossible to design at this specific speed, but still need to explain lock's assertion that "The efficiencies of forced vortex pumps in the same specific speed range and capacity are higher than those of similar conventional pump"}

Barske comments that the number of impeller blades does not seem to be critical, with 3-6 being operated successfully. However, altering the number of blades can be used to elminite fluid vibrations \cite{barske1953}. Lock provides corrections for blades from 1 to 16 \cite{lockfig10A}.

The open impeller choice is to further reduce disc friction losses \todo{is this true?} by minimizing surface area. Conventional impellers uses shrouds to guide the flow and reduce secondary flows. However, as in a Barske pump the impeller serves only to generate the forced vortex flow, this was deemed unnecessary by at least Barske and Lock. This also removes axial thrust loads on the impeller, helping with sealing and mechanical design of the shaft. Also, simplify manufacturing.

Conventional pump design with current design tools produce nonsensical geometries and results at low specific speeds because of one number that is used in the design process, which is the slip coefficient. The slip coefficient determines the angular deflection of the fluid leaving the impeller from the blade angle at the impeller exit, because a finite number of blades cannot perfectly guide the flow. It is given by an empirical regression, the most widely used being Wiesner's fit to Busemann's potential flow solutions \cite{gulichCentrifugalPumps2020}:

\begin{equation}
    \sigma = 1 - \frac{\sqrt{\sin\beta_{2B}}}{z^{0.7}}
    \label{eq:wiesner_slip}
\end{equation}

where $\beta_{2B}$ is the blade exit angle from the circumferential direction and $z$ the blade count. The slip factor closes the theoretical head, because the Euler head of \cref{eq:euler_specific_work} with a finite blade count and the outlet velocity triangle becomes

\begin{equation}
    \psi_{th} = 2\left(\sigma - \phi_2 \cot\beta_{2B}\right)
    \label{eq:head_with_slip}
\end{equation}

with $\psi$ and $\phi_2$ the head and flow coefficients of \cref{eq:psi_phi}. The difficulty at low specific speed is now visible in the equation rather than asserted. The duty fixes a high target $\psi$ at a very low $\phi_2$, which \cref{eq:head_with_slip} can only meet by pushing $\beta_{2B}$ towards the radial and the blade count up, while continuity $\phi_2 = Q/(\pi d_2 b_2 u_2)$ at the same low flow and large diameter simultaneously drives the outlet width $b_2$ towards zero. The closure does have a root, but it lands on geometry that is either unmanufacturable or outside the validity range of \cref{eq:wiesner_slip}, whose own limit on the blade loading is $\varepsilon_{lim} = \exp(-8.16\sin\beta_{2B}/z)$. This is what motivates the forced vortex Barske architecture, which abandons the attached blade-guided flow model entirely and so removes the slip factor from the design process.

This claim needs to be scoped, because a formally converged conventional design at this specific speed has been published. Kim obtained a converged Wiesner design with $\beta_{2B} = 51.7^\circ$ at the same specific speed as this thesis, but only by mixing correlations from different sources, each used outside its stated validity range, and the resulting geometry required closed channels of $b_2 = 0.31$ mm on a 24 mm impeller \cite{kimDevelopmentUltraLowSpecific2022}. That is exactly the corner of \cref{eq:head_with_slip} described above: a root exists, but the practical failure of the conventional architecture at this scale is manufacturability and suction robustness rather than pure mathematics. This is explored further with the design attempt in \cref{sec:appendix_conventional} and \cref{sec:results}.

\subsubsection{Euler and Forced-Vortex Head}

\begin{figure}[H]
    \centering
    
\includegraphics[width=0.55\textwidth]{placeholder.png}
    \caption{Velocity triangle nomenclature and the forced vortex flow field in the Barske impeller, where the fluid rotates as a solid body with $v_\theta = \omega r$.} % PLACEHOLDER IMAGE
    \label{fig:pump-velocity-triangle}
\end{figure}

All theories relating to pump and turbine design can be derived from the conservation of angular momentum by Newton's second law of motion. In turbomachinery, this is usually expressed in terms of the Euler turbine equation:

\begin{equation}
    P=\tau\omega=\rho Q \left(r_2 c_{2u} - r_1 c_{1u}\right)\omega=\rho Q\left(u_2c_{2u}-u_1c_{1u}\right)
    \label{eq:euler_turbine_equation}
\end{equation}

Rewriting this equation in terms of specific work $Y=P/\rho Q$ per unit mass:

\begin{equation}
    Y = \left(r_2 c_{2u} - r_1 c_{1u}\right)\omega = u_2c_{2u}-u_1c_{1u}
     \label{eq:euler_specific_work}
\end{equation}

derive from bernouli, with rotating frame substitution $c^2=w^2-u^2+2\times u \times c_u$ and neglecting compressibility losses:

$$p_{1}+\frac{\rho}{2}w_{1}^2-\frac{\rho}{2}u_{1}^2=p_{2}+\frac{\rho}{2}w_{2}^2-\frac{\rho}{2}u_{2}^2$$

and as such the theoretical rise in static pressure:

\begin{equation}
    \Delta p_{st} = p_2-p_1 = \frac{\rho}{2}\left(w_1^2-w_2^2\right)+\frac{\rho}{2}\left(u_2^2-u_1^2\right)
    \label{eq:pressure_rise}
\end{equation}

In Barske, Lobanoff and Lock the relative velocity $w$ is assumed to be negligable, ending up with:

\begin{equation}
    \Delta p_{st} = \frac{\rho}{2}\left(u_2^2-u_1^2\right)
    \label{eq:pressure_rise_negligible_w}
\end{equation}

Then, the kinetic energy increase assuming negligible inlet prerotation is:

\begin{equation}
    \Delta p_{st} = \frac{\rho}{2}u_2^2
    \label{eq:pressure_rise_negligible_w_prerotation}
\end{equation}




\cref{eq:pressure_rise_negligible_w_prerotation} is the theoretical total pressure a barske impeller generates at the blade tips. In a barske pump, the diffuser is responsible for converting this kinetic energy into static pressure by slowing down the fluid in an expanding diffuser. Barske introduces a loss coefficient for the kinetic energy which can be recovered through the diffuser to be between 0.2 and 0.5. Barske uses the symbol $\psi$ for this coefficient, but as it conflicts with the commonly used symbol for head coefficient, this paper will instead use $k_d$. Combining the static and kinetic pressure rise, the total increase in pressure is:

\begin{equation}
    \Delta p = \Delta p_{st} + k_d\, q=\frac{1}{2}\rho\left((1+k_d) u_2^2 - u_1^2\right)
    \label{eq:barske_pressure_rise}
\end{equation}

This forms a simple basic model for barske pump performance. 

\subsubsection{Barske Sizing Method}
\label{sec:theory_barske_sizing}

The barske pump is sized by the methods provided by barske \cite{barske1953} with the assitance of Gulich. First the inlet is sized by choosing a reasonable inlet velocity and calculating the inlet area:

\begin{equation}
    d_0 = \sqrt{\frac{4\dot{m}}{\pi \rho c_0}}
    \label{eq:barske_inlet_diameter}
\end{equation}

The impeller blade inlet diameter is enlarged by a small amount as some pre-rotation is desirable for increased head. For example, $d_1\sim1.1d_0$.

By substituting $u=\pi n d / 60$ into \cref{eq:barske_pressure_rise}, the required impeller tip diameter can be calculated:

\begin{equation}
    d_2 = \sqrt{\frac{1}{1+k_d}\left(\frac{7200\Delta p}{\rho n^2\pi^2} + d_1^2\right)}
    \label{eq:barske_tip_diameter}
\end{equation}

Barske then gives these guidelines for other dimensions of the pump:

\begin{align}
    s_{ax} & \leq \min{\left(0.75, \frac{d_2}{100}\right)} \\
    b_1 & \geq \frac{1}{4} d_1 \\ %note to self, b is blade height
    b_2 & \geq \max{\left(b_1 \frac{d_1}{d_2}, 3s_{ax}\right)} \\
\end{align}

For the annular casing sizing, Barske doesn't have very specific numbers, so Gulich's guideline where B/H = 1.5 to 2.5 can be followed:

\begin{align}
    B & = 2 s_ax + b_2 \\
    H & = B / 2
\end{align}

\todo{look at efficiency that dynamicpumps uses, and cavitation maybe as well, iirc they have first principles}

To size the diffuser, Barske suggests sizing the diffuser throat to be around 0.8 to 0.9 of $u_2$, with the area of the diffuser outlet to be 3-4 times the area of the diffuser throat.

\subsubsection{Efficiency and Disc Friction}

Barske assumes that the disc friction losses can be approximated as a circular disc in a rotating fluid. Converting the equation into metric units gives:

\begin{equation}
    P_{RR} [\mathrm{W}] = 1956\times\rho\nu^{0.2}\left(\frac{n}{1000}\right)^{2.8}\left(d_2^{4.6}+4.6d_1^{3.6}b_1\right)
\end{equation}

As the hydraulic efficiency is defined as the useful power out of the pump over the total power into the pump, efficiency can be defined as:

\begin{align}
\eta_h & = \frac{\text{Useful power}}{\text{Input power}} = \frac{P}{P_{th}+P_{RR}}
 & = \frac{Q\Delta p}{Q\Delta p_{th}+P_{RR}}
 & = \frac{\Delta p}{\Delta p_{th}+P_{RR}/Q}
\end{align}

If total shaft efficiency is desired, the mechanical losses can be included:

\begin{equation}
    \eta=\frac{\Delta p}{\Delta p_{th}+\frac{P_{RR}+P_m}{Q}}
\end{equation}

\subsubsection{Lock's Refined Model}

Lock provides a refined model which takes into account prerotation, nozzle losses and correction coefficients taking into account theoretical two dimensional flow solutions by A. Busemann \cite{buseman}. \todo{think if his 5 or so equations should be included here ot not, he also says that disc friction is treated as mechanical loss so doesn't effect head measurements and only effect efficiency calcluations, i will check if my code does this later on.}


By writing the total theoretical pressure gain in head:

\begin{equation}
    H_{th,0} = \frac{u_2^2 - u_1^2}{2g} + \frac{u_2^2}{2g}
\end{equation}

Lock uses busemann's corrections and derives a refined total head, taking into account prerotation:

\begin{equation}
    H_{th} = k_0\frac{u_2^2}{g} - C_H\,K_1\!\left(1-\frac{Q}{Q_{opt}}\right)\!\frac{d_1}{d_2}\,\frac{u_2^2}{g}
    \label{eq:lock_total_head}
\end{equation}

\begin{equation}
H = H_{th} - \underbrace{\zeta_{D}\left[1-\left(\tfrac{d_3}{d_4}\right)^2\right]^2\frac{c_3^2}{2g}}_{\text{diffuser-nozzle loss}} - \frac{c_4^2}{2g} + \frac{c_0^2}{2g}
\label{eq:lock_static_head}
\end{equation}

Where $\zeta_{D}=0.194$ is the diffuser loss coefficient from Hydraulic Institute's Pipe Friction Manual \cite{hydraulicpipemanual}.

By expressing $c^2/2g$ by $(8/g\pi^2)Q^2/d^4$ with $c=4Q/\pi d^2$, the head can be expressed as a function of flow rate:

\begin{equation}
    H = k_0\frac{u_2^2}{g} - C_H\,K_1\!\left(1-\frac{Q}{Q_{opt}}\right)\!\frac{d_1}{d_2}\,\frac{u_2^2}{g} - \frac{8Q^2}{g\pi^2}\left(\frac{\zeta_{D}}{d_3^4} + \frac{1}{d_4^4} - \frac{1}{d_0^4}\right)
    \label{eq:lock_static_head_flow}
\end{equation}

To find the optimum flow rate used to shape the pump curve, the derivative of the head with respect to flow rate can be taken and set to zero:

\begin{equation}
    Q_{opt} = \sqrt{\frac{k_0\frac{u_2^2}{g}}{\frac{24}{g\pi^2}\left(\frac{\zeta_{D}}{d_3^4}\left(1-\left(\tfrac{d_3}{d_4}\right)^2\right) + \frac{1}{d_4^4} - \frac{1}{d_0^4}\right)}}
    \label{eq:lock_optimum_flow}
\end{equation}

Lock and Barske suggest that the flow cutoff point is limited by the diffuser throat. To find the static pressure at the diffuser throat, reversing the transfer between dynamic head into static head provided by the previous equations \todo{link to the previous equations}:

\begin{equation}
    H_3 = H + \frac{8Q^2}{g\pi^2}\left(\frac{\zeta_{D}}{d_3^4}\left(1-\left(\frac{d_3}{d_4}\right)^2\right) + \frac{1}{d_4^4} - \frac{1}{d_3^4}\right)
    \label{eq:lock_diffuser_throat}
\end{equation}

Then solving for the $Q$ at which $H_3$ is equal to vapor pressure head $H_{vap}=\frac{1}{\rho g}\left(p_{vap}-p_{inlet}\right)$ gives the flow cutoff point.

This cutoff mechanism is the key difference between the two pump models used in this thesis. The Euler and Barske models of \cref{eq:pressure_rise_negligible_w_prerotation} predict a flat head curve with no flow dependence at all. Lock's model predicts a drooping head curve, and on top of that it predicts that the curve ends at a maximum flow rate where the static head at the diffuser throat reaches vapour pressure and the throat cavitates. Both predictions are directly testable, and the comparison of the measured head curve against both models is one of the main results of \cref{sec:results}. Two limitations that Lock states himself should be kept in mind when his model is compared to data later. He notes that rounding of the blade leading edges reduces the effective head, and that disc friction is treated as a mechanical loss, so it does not affect the head equations and only enters the efficiency \cite{lockForcedVortexPump1966}. These acknowledged but unmodelled effects matter when the model constants are fitted to data in \cref{sec:results}.


\subsubsection{Non-dimensional Performance and Scaling}
\label{sec:theory_nondim}

The measured performance in \cref{sec:results} is presented in non-dimensional form, so the coefficients are defined here. Following Gülich \cite{gulichCentrifugalPumps2020}, the head and flow coefficients are

\begin{equation}
    \psi = \frac{2gH}{u_2^2}, \qquad
    \phi = \frac{Q}{\pi d_2 b_2 u_2}
    \label{eq:psi_phi}
\end{equation}

The head coefficient measures how much of the natural pressure scale of the rotor, the dynamic head of its tip speed, the machine converts. It is particularly meaningful for this pump: a pure forced vortex delivers $\Delta p = \tfrac{1}{2}\rho u_2^2$, which is exactly $\psi = 1$, so the measured $\psi$ directly tests whether the forced vortex is achieved. The flow coefficient fixes the shape of the velocity triangles, so two operating points with the same $\phi$ have the same flow physics regardless of speed.

For a fixed geometry, dimensional analysis of $(Q, H, n, d_2, \rho, \mu)$ reduces the pump behaviour to $\psi = f(\phi, Re)$ with the tip Reynolds number

\begin{equation}
    Re_u = \frac{u_2\, r_2}{\nu}
    \label{eq:re_u}
\end{equation}

Because pump head is Reynolds insensitive to first order, all speeds should collapse onto a single $\psi(\phi)$ curve. Holding $\psi$ and $\phi$ constant between two speeds gives the affinity laws,

\begin{equation}
    Q \propto n, \qquad H \propto n^2, \qquad P \propto n^3
    \label{eq:affinity}
\end{equation}

which are the licence for any extrapolation across speed. The losses are not Reynolds insensitive, which is why the efficiency does not collapse the way the head does, and the loss scaling is treated empirically through a loss ratio of the form $(1-\eta) \propto Re^{-a}$.

For cavitation, the suction state is expressed as the net positive suction head,

\begin{equation}
    \textrm{NPSH}_a = \frac{p_{in} - p_{vap}}{\rho g} + \frac{v_{in}^2}{2g}
    \label{eq:npsha}
\end{equation}

the absolute margin of the inlet head above vapour pressure. The standard performance criterion is NPSH$_3$, the available NPSH at which the developed head has dropped 3\% below its fully wetted value at constant speed and flow \cite{gulichCentrifugalPumps2020}. Suction capability is compared across machines through the suction specific speed, defined like \cref{eq:specific_speed} but with NPSH$_3$ in place of head.

\subsection{Turbine Theory}

\subsubsection{Impulse Turbines}

\begin{figure}[H]
    \centering
    \includegraphics[width=0.6\textwidth]{turbine_cantilever.pdf}
    \caption{Impulse turbine velocity triangle nomenclature at the rotor inlet (station 3) and exit (station 4). The design point triangles with actual values are shown in \cref{fig:turbine-triangles-design}.} % PLACEHOLDER IMAGE: generic labelled triangles, or merge with fig:turbine-triangles-design
    \label{fig:turbine-velocity-triangles}
\end{figure}

Turbines work with the same Euler turbine equation. In the supersonic impulse turbine design, the blades are only occupy a small outer portion of the blisk, so the turbine equation can be simplified:

\begin{equation}
    \Delta h_{\textrm{use}}=c_{3u}u_3-c_{4u}u_4=u(c_{3u}-c_{4u})
    \label{eq:turbine_specific_work}
\end{equation}

From \cref{eq:turbine_specific_work}, it is understood that a higher flow velocity difference $c_{3u}-c_{4u}$ per unit mass will lead to a higher $\Delta k_0$ according to equation \cref{eq:euler_specific_work}. This means that high absolute velocities, and for the maximum amount of the flow angled in the circumferential direction is optimal. For this reason, gas generator cycle turbines tend to take advantage of the high pressure ratios and thus the supersonic velocities. Then for the angles, the highest turning angle achievable is optimal. The amount of flow turning is usually limited by the achievable geometry.

The ideal efficiency of the impulse stage follows from \cref{eq:turbine_specific_work} and the velocity triangles. At the rotor inlet and exit the circumferential components of the absolute and relative velocities are related by

\begin{equation}
    c_{3u} = w_{3u} + u, \qquad c_{4u} = u + w_{4u}
    \label{eq:triangle_components}
\end{equation}

The ideal impulse blade turns the relative flow without slowing it down, so the exit relative velocity mirrors the inlet one, $w_{4u} = -w_{3u}$. Substituting into \cref{eq:turbine_specific_work},

\begin{equation}
    \Delta h_{\textrm{use}} = u\left(c_{3u} - c_{4u}\right)
    = u\left[(w_{3u}+u) - (u+w_{4u})\right]
    = 2u\, w_{3u}
    = 2u\left(c_{3u} - u\right)
    \label{eq:impulse_work}
\end{equation}

Two things come out of \cref{eq:impulse_work}. First, solving it for the inlet swirl gives the velocity triangle inversion used by the sizing procedure of \cref{sec:design}: given a required specific work and a blade speed, the nozzle must deliver

\begin{equation}
    c_{3u} = \frac{\Delta h_{\textrm{use}}}{2u} + u, \qquad
    c_3 = \frac{c_{3u}}{\cos\beta_3}, \qquad
    c_{3m} = c_3 \sin\beta_3
    \label{eq:triangle_inversion}
\end{equation}

where $\beta_3$ is the nozzle angle measured from the circumferential direction. Second, dividing the work by the kinetic energy the nozzle supplies gives the stage efficiency. In a pure impulse stage the entire expansion happens in the nozzle, so the available energy per unit mass is the jet kinetic energy $c_3^2/2$, and with $c_{3u} = c_3\cos\beta_3$, \todo{this is wrong. wtf is inlet swirl}

\begin{equation}
    \eta_{\textrm{ideal}}
    = \frac{2u\left(c_3\cos\beta_3 - u\right)}{c_3^2/2}
    = 4\,\frac{u}{c_3}\left(\cos\beta_3 - \frac{u}{c_3}\right)
    \label{eq:impulse_ideal_efficiency}
\end{equation}

The efficiency depends only on the blade to jet speed ratio $u/c_3$ and the nozzle angle. By taking the derivative of \cref{eq:impulse_ideal_efficiency} and setting it to zero, the optimum sits at $u/c_3 = \cos\beta_3/2$, where the exit swirl is zero and the efficiency peaks at $\cos^2\beta_3$. The same expression with the Weiss loss coefficients of \cref{subsec:weissloss} attached becomes $\eta_h = 2k_N^2\,(u/c_3)\left(\cos\beta_3 - u/c_3\right)(1+k_R)$, which reduces to \cref{eq:impulse_ideal_efficiency} when the nozzle and rotor are loss free ($k_N = k_R = 1$); this lossy form is the solid curve in \cref{fig:turbine-eta-uc}.

\begin{figure}[H]
    \centering
    \includegraphics[width=0.6\textwidth]{turbine_theory.pdf}
    \caption{Impulse stage efficiency against blade to jet speed ratio. The dashed line is the ideal impulse stage of \cref{eq:impulse_ideal_efficiency}, the solid line includes the Weiss loss coefficients evaluated at the design point, and the star marks the design point of this turbine, far below the optimum speed ratio.}
    \label{fig:turbine-eta-uc}
\end{figure}


% In my notes, this can be found "SP-8110.md" file and "Supersonic Turbine Blade Pitch and Height.md"
The jet velocity $c_s$ used in this ratio is the spouting velocity, the velocity reached by an ideal isentropic expansion of the gas over the full stage pressure ratio,

\begin{equation}
    c_s = \sqrt{2\,\Delta h_s}
    = \sqrt{2\, c_p T_{01} \left[1 - \left(\frac{p_e}{p_{01}}\right)^{\frac{\gamma-1}{\gamma}}\right]}
    \label{eq:spouting_velocity}
\end{equation}

It links the available inlet pressure directly to the achievable jet velocity, and it is what makes the blade speed problem concrete. A small turbine on a 20,000 RPM shaft has a blade speed of around 100 m/s, while the spouting velocity of even a moderate pressure ratio gas generator is several hundred m/s, so $u/c_s$ is forced far below the optimum of around 0.5.

NASA special publication 8110 shows that for low $u/c_s$ turbines, impulse turbines are more efficient, even if reaction turbines can have a higher overall efficiency, just not at the low $u/c_s$ regime. In addition, the same conclusion is drawn by comparing the influence of tip clearance on the two types, where impulse turbines are more robust to tip clearance. This is especially important for small scale turbines where manufacturing tolerances are more difficult to achieve and contribute to more significant losses \cite{NASASP8110}. \todo{need to explain what reaction and impulse turbines are, maybe with a diagram?}



\subsubsection{Blade Height and Partial Admission}

Another problem with designing small turbines is with the low blade height required. By calculating the mass flux through the turbine, the blade height can be designed:

\begin{equation}
    H_3 = \frac{\dot{m}}{\rho_3 c_{3m} d_{\textrm{mean}}}
    \label{eq:turbine_blade_height_full_admission}
\end{equation}

Where $\rho_3$ is the density of the turbine gas at the turbine inlet. Using the example turbine above, the required blade height would be \todo{addvalue}. By introducing partial admission, the mass flow through the turbine can be restricted to a smaller arc around the annulus at a slight cost to efficiency. The partial admission ratio $\zeta$ can be included into the blade height formula as so:

\begin{equation}
    H_3 = \frac{\dot{m}}{\zeta \rho_3 c_{3m} d_{\textrm{mean}}}
    \label{eq:turbine_blade_height}
\end{equation}

By reducing the partial admission to \todo{addvalue}, the blade height is increased to \todo{addvalue}. The design of the turbine blade height is very important due to multiple reasons. Firstly, the aspect ratio of the turbine blade determines the magnitude of of the wall boundary layer blockage. For this reason, a aspect ratio between \todo{addvalue} is desired.

Another way to increase the blade height is to decrease the meridonal velocity $c_{3m}$. This is achieved by increasing the angle that the flow has to turn as $c_{3m}=c_3\sin(\beta_3)$.

The cost of partial admission and of low aspect ratio blading is not guesswork, as both have experimentally validated loss theories. Ohlsson derived and validated a closed form partial admission loss model where the efficiency penalty is a function of the ratio of active arc length to blade chord \cite{ohlssonPartialAdmissionTurbines1962}, and the same rig produced loss correlations for very low aspect ratio impulse blades \cite{ohlssonLowAspectRatio1964}. These are used to check the loss model predictions in \cref{sec:results}.

For supersonic flow to develop within a turbine blade passage, a condition called "starting" must be satisfied. On the startup of a supersonic turbine, a normal shock starts at the inlet of the blade passage and as to pass through the whole passage before the flow into and out of the turbine passage can be supersonic. \todo{add derivation for this, later? also, loss models.}


\subsubsection{Loss Model}
\label{subsec:weissloss}

The main loss model used in the design of the turbine is the one provided by Weiß \cite{weiss}. There are more general loss models provided by texts such as NASA SP-290, giving loss models for tip-clearance, disk friction, partial admission and incidence losses. However, those models were built on large full-admission, subsonic to transonic turbines. Weiss specifically derived his loss model for the specialized application of a supersonic impulse turbines at small scales (1-200kW). As the regime of the turbine considered in this thesis matches closer to Weiss's model compared to SP-290's model, Weiss's model is used. Weiss describes different models lead to stage efficiency differences of about 10\% points. Derived from rocket turbopumps literature and adjusted by in-house measurements, Weiss provides three efficiencies.


\begin{equation}
    k_N = \sqrt{1 - \left( 0.0029\,M_3^3 - 0.0502\,M_3^2 + 0.2241\,M_3 - 0.0877 \right)}
    \label{eq:weiss_nozzle_coeff}
\end{equation}

\cref{eq:weiss_nozzle_coeff} is the nozzle loss coefficient, relating the actual velocity to the theoretical isentropic velocity by $k_N=c_{No}/c_{No,is}$. 

\begin{equation}
    \begin{aligned}
        k_R ={}& 0.957
            - 3.62\times10^{-4}\,\Delta\beta
            - 2.58\times10^{-2}\,M_{w3} \\
            &+ 6.39\times10^{-6}\,\Delta\beta^2
            + 6.74\times10^{-2}\,M_{w3}^2 \\
            &- 7.53\times10^{-8}\,\Delta\beta^3
            - 4.30\times10^{-2}\,M_{w3}^3 \\
            &- 2.38\times10^{-4}\,\Delta\beta\, M_{w3}
            + 1.45\times10^{-6}\,\Delta\beta^2 M_{w3}
            + 4.25\times10^{-5}\,\Delta\beta\, M_{w3}^2
    \end{aligned}
    \label{eq:weiss_rotor_coeff}
\end{equation}

\cref{eq:weiss_rotor_coeff} is the loss of velocity along the blade passage through a theoretically impulse blade, relating the actual velocity to the theoretical velocity by $k_R=w_{exit}/w_{inlet}$.

\begin{equation}
    P_v = \frac{1.85}{2}\,(1 - \varepsilon)\,\rho_3 \left( \frac{n}{60} \right)^{3} d_m^4 \,(4.5\,b)
    \label{eq:weiss_ventilation}
\end{equation}

Finally, \cref{eq:weiss_ventilation} captures the losses due to partial admission. Weiss calls this loss ventillation losses, but it can also be called windage losses. Ventilation losses refers to the the non-admitted arc being dragged through stagnant air.

The choice of Weiss's model over the classical decompositions is a regime match rather than a claim that one is better. The losses that dominate a small supersonic partial admission impulse stage are supersonic friction in the nozzle and blade passage plus ventilation, which is exactly what Weiss's coefficients parameterise. The classical loss breakdowns such as Glassman's are anchored in large full admission machines where tip clearance and incidence dominate \cite{glassmanTurbineDesignApplication1994}, and those losses are largely absent here by architecture, since the blades are shrouded and the velocity triangles are matched at the design point. Whether the Weiss prediction is actually trustworthy is checked against measured data from a machine of the same architecture in \cref{sec:results}.

\todo{FUCKS SAKE LMAO, look at comment}
% Weiss says this: There are somewell-known turbine loss models in literaturee.g. by Refs. [39e41] etc. All of them are mainly developed for and based on the data from subsonic MW-class multistage steam and gas turbines working with full admission. 
% 39 - Traupel W Thermische, 40 - Craig HRM, 41 - Kacker SC


\subsubsection{Method of Characteristics Blade Design}
\label{subsec:method_of_characteristics}

The design so far is based on 1D mean line theory. The design of the actual blade profile which will turn the supersonic flow in the desired direction is often accomplished using an application of the Method of Characteristics to transition a uniform flow field into vortex flow. \todo{add some history here, iirc sp290 has some stuff on why vortex flow, but here needs some stuff about how vortex flow was proven to be able to turn flows, and method of characteristics is used to turn the flow into such vortex flow, so obviously some explanation of vortex flow is required}. 

\paragraph{Blade terminology and vortex flow}

Geometry is separated into the turbine blade, and the passage between the flow. Two surfaces are separately designed, and it is common to refer to the top surface as the convex, and the bottom surface as the concave one. This is in relation to the solid blades themselves as opposed to an alternate designation by the top and bottom of the blade passage.

Vortex flow is a special type of supersonic flow in which flow follows a theoretical perfect arc around a central point. This is a very special type of flow field that is achieved by conditions where expansion and compression waves cancel each other to allow shockless rotation of the flow field. The vortex flow is bounded by the two circular arcs of the surfaces of which a desired surface Mach number is achieved. MOC generates the profile for desired mach numbers. Later analysis will need to be required to choose these desired Mach numbers.

A special condition of vortex flow satisfies the equation

\begin{equation}
    \frac{c}{r} = \textrm{constant}
    \label{eq:vortex_condition}
\end{equation}

Boxer and Goldman nondimensionalizes this equation using the critical velocity $M^{*}=c/a^*$ where $a^*=\sqrt{\gamma R T^*}$ is the speed of sound at the sonic point. As $a^*$ only depends on stagnation temperature, it is constant along any streamline of fixed stagnation enthalpy. Then, a nondimensional radius $r^*$ is introduced as the sonic velocity streamline in the vortex field, giving the following condition:

\begin{equation}
    \frac{M^*}{r^*} = 1
    \label{eq:vortex_condition_nondim}
\end{equation}

The Prantdl-Meyer angle can be defined from the critical mach number $M^*$ as:

\begin{equation}
\nu(M) = \frac{\pi}{4}\left(\sqrt{\frac{\gamma+1}{\gamma-1}} - 1\right) + \frac{1}{2}\left(\sqrt{\frac{\gamma+1}{\gamma-1}} \arcsin\left[(\gamma-1)M^{*2} - \gamma\right] + \arcsin\left(\frac{\gamma+1}{M^{*2}} - \gamma\right)\right)
\label{eq:prandtl_meyer_angle}
\end{equation}

And the converting between Mach number and Critical Mach number is given by:

\begin{equation}
M^* = \sqrt{\frac{(\gamma+1)M^2/2}{1 + (\gamma-1)M^2/2}}
\label{eq:mach_to_critical_mach}
\end{equation}

Equations \cref{eq:vortex_condition_nondim} and \cref{eq:mach_to_critical_mach} can be used to derive the nondimensional radius of the upper and lower surfaces of the blade. The end points of this circular arc is exactly the angle by which the inlet flow has to be turned to achieve the surface mach number. 

\begin{equation}
\begin{aligned}
    \alpha_{l-i} & = \beta_1 - (\nu_i - \nu_l) \\
    \alpha_{u-i} & = \beta_1 + (\nu_i - \nu_u)
\end{aligned}
\label{eq:inlet_flow_turning_angle}
\end{equation}

The equations for the outlet are the same as in \cref{eq:inlet_flow_turning_angle}.

\paragraph{Transition arcs}

To transition from the uniform inlet flow to the vortex flow, method of characteristics is used to generate the blade profile to achieve this in a shockless manner. 

\todo{explain new axis}

To achieve this, the flow angles from $\beta_i$ to $\alpha_i$ is discretized into $k$ interval of $\phi_0$ to $\phi_k$. It is shown by \todo{find the ref from goldman 1968 ref 2} that the flow angle and dimensionless radius $R^*$ is related along its characteristic line with:

\begin{equation}
    \phi = \pm \frac{1}{2}f(R^*) + \textrm{constant}
    \label{eq:characteristic_line_equation}
\end{equation}

The positive gives the expansion line and the negative gives the compression line. Where $f(R^*)$ is:

\begin{equation}
    f(R^*) = \sqrt{\frac{\gamma + 1}{\gamma - 1}}\arcsin\left(\frac{\gamma-1}{R^{*2}} - \gamma\right) + \arcsin\left(\left(\gamma+1\right)R^{*2} - \gamma\right)
    \label{eq:f_R_star}
\end{equation}

$f(R^*)$ for a given flow angle can be calculated with:

\begin{equation}
    f(R^*_k) = 2\nu_i - \frac{\pi}{2}\left(\sqrt{\frac{\gamma+1}{\gamma-1}} - 1\right) - 2k\Delta\nu
    \label{eq:f_R_star_k}
\end{equation}

The result of \cref{eq:f_R_star_k} can be used to numerically solve for $R^*_k$ with \cref{eq:f_R_star}, and points corresponding to each flow angle along the major expansion characteristic line can be calculated with:
\begin{equation}
    \begin{aligned}
    x^*_{k} & = -R^*_k \sin(\phi_k) \\
    y^*_{k} & = R^*_k \cos(\phi_k)
    \end{aligned}
\end{equation}

The wall segments are generated by starting from the end point of the circular arc, then finding the intersection of the previous wall segment and the mach line from this point. The slope of the mach line with respect to to the $x^*$ axis is given by:

\begin{equation}
    m_k = \tan\left(\frac{\phi_k+\phi_{k-1}}{2}+\frac{\mu_k+\mu_{k-1}}{2}\right)
\end{equation}

Where
\begin{equation}
    \mu_k = -\arcsin\left(\frac{1}{M_k}\right) = -\arcsin\left(\sqrt{\left(\frac{\gamma+1}{2}\right)R^{*2}_k - \frac{\gamma-1}{2}}\right)
\end{equation}

By iterating through the flow angles for both the upper and lower surfaces for the inlet and outlet, the full theoretical MOC blade profile can be generated.

\todo{add "i" to the variables to define inlet}

\subsubsection{Surface Mach Number Selection}

There are two main factors to consider when choosing the inlet, outlet and surface mach numbers. One is to prevent flow separation, and the other is to ensure that supersonic flow within the blade passage can "start".

\paragraph{Flow separation}

Separation occurs in high adverse pressure gradients. In supersonic flow, this happens whenever the flow is decelerating along the surface.  There are two main regions on the blade surfaces which separation can occur. At the upper surface, the inlet transition arc accelerates the flow from the inlet velocity until the vortex flow, then decelerates the flow until the outlet. At the lower surface, the inlet transition arc decelerates the flow from the inlet velocity until the vortex flow, then accelerates the flow until the outlet. 

\todo{add a figure, question: literature figure or my figure? probably both is best}

Goldman believes that separation on the lower surface is probably not as important as the flow would tend to reattach shortly as the pressure gradient becomes favourable again \cite{goldman1970}. Goldman uses the value of the incompressible form factor $H_i$ along the blade surface to predict separation, where  Schlichting found \todo{find goldman ref 12} that separation usually occurs in the range of $H_i$ from 1.8 to 2.4.

To find $H_i$ along both surfaces, Sasman and Cresci \cite{sasmancresci} provides a method to calculate the compressible turbulent integral boundary layer with two ODEs. 

Once a MOC blade profile is generated and the mach number at every point along the blade is known, the momentum-integral equation and shape-factor equation can be solved as an initial value problem starting from the leading edge goverened by the following equations:

\begin{equation}
    \frac{df}{d(s/c)} = 1.268\left[-\frac{f}{M_e}\frac{dM_e}{d(s/c)}\left(1 + g_w H_i\right) + A\,c\right]
\end{equation}

\begin{equation}
    A\,c = 0.123\,e^{-1.561 H_i}\, M_e\, \mathrm{Re}_0\, \frac{T_e}{T_0}\left(\frac{\bar\mu}{\mu_0}\right)^{0.268}
\end{equation}

\begin{equation}
\frac{dH_i}{d(s/c)} = \underbrace{-\frac{1}{2M_e}\frac{dM_e}{d(s/c)}\,H_i(H_i+1)^2(H_i-1)\left[1 + (g_w-1)\frac{H_i^2+4H_i-1}{(H_i+1)(H_i+3)}\right]}_{\text{pressure-gradient term: deceleration drives } H_i \text{ up}} \\ + \underbrace{\frac{H_i^2-1}{f}\,A\,c\left[H_i - \frac{0.011\,(H_i+1)(H_i-1)^2}{H_i^2\,(C_f/2)}\frac{T_e}{\bar T}\right]}_{\text{shear term: wall friction relaxes } H_i \text{ back toward equilibrium}}
\end{equation}

With initial conditions of a flat-plate $\frac{\theta_0}{c} = 0.036\,\frac{s_0}{c}\,\mathrm{Re}_{s_0}^{-1/5}$. \todo{cite this}
% i found delta = 0.37x/(Re to 1/5) from anderson 19.1 page 1052, prandtl u/U = (y/delta) to 1/7 then shoving into the equation of momentum thickness gives 7/72 delta, which is 0.0359722222, done.

The initial condition for the shape factor is chosen so that it matches closely to the rapidly converging value $ H_{i,0} = 1.67$. \todo{make this more rigorous}

It is important to note that Goldman provides equations in his NASA TN D-4421 1968 paper that has been shown to not be conservative enough, with tests by Erikson saying that his criteria is a little too optimistic. The method described here is from Goldman's NASA TM X-2095 paper.

Two limitations of the method should be stated, both flagged by Goldman himself \cite{goldmanAnalyticalInvestigationBlade1970}. The $H_i$ threshold band of 1.8 to 2.4 is rooted in incompressible data, and the method assumes a turbulent boundary layer at a chord Reynolds number of $3.5\times10^4$ in his reference case, where laminar flow is plausible. The flight programme experience supports running the check on both surfaces, as the HM60 supersonic stage found boundary layer growth considerable with separation starting at around 25\% chord on the convex side, which forced a rotor redesign \cite{erikssonDevelopmentSupersonicTurbine}.

\todo{reorganize here}

\todo{should I remove these?}
%  it took me a while to understand, but I think I finally understand it. Sasman and cresci developed a way to transform (stretch) the surface arc coordinates from compressible to incompressible coordinates, somehow deriving their eq 18, 19, 20 which can be solved as an IVP with initial conditions in 

\begin{figure}[H]
    \centering
    \begin{subfigure}[b]{0.49\textwidth}
        \centering
        \includegraphics[width=\textwidth]{goldman_paper_fig4.png}
        \caption{Goldman's published distribution \cite{goldmanAnalyticalInvestigationBlade1970}}
        \label{fig:goldman-mach-paper}
    \end{subfigure}
    \hfill
    \begin{subfigure}[b]{0.49\textwidth}
        \centering
        \includegraphics[width=\textwidth]{goldman_validation_mach.pdf}
        \caption{This work's implementation}
        \label{fig:goldman-mach-ours}
    \end{subfigure}
    \caption{Validation of the MoC implementation against Goldman's typical design case ($M_{in}=2.5$, $\beta=70^\circ$). The surface Mach plateaus reproduce to within 1\% (2.96 and 1.84 against Goldman's 2.97 and 1.85).}
    \label{fig:goldman-validation-mach}
\end{figure}

\begin{figure}[H]
    \centering
    \begin{subfigure}[b]{0.49\textwidth}
        \centering
        \includegraphics[width=\textwidth]{goldman_paper_fig5.png}
        \caption{Goldman's published form factor \cite{goldmanAnalyticalInvestigationBlade1970}}
        \label{fig:goldman-hi-paper}
    \end{subfigure}
    \hfill
    \begin{subfigure}[b]{0.49\textwidth}
        \centering
        \includegraphics[width=\textwidth]{goldman_validation_hi.pdf}
        \caption{This work's implementation}
        \label{fig:goldman-hi-ours}
    \end{subfigure}
    \caption{Validation of the Sasman and Cresci boundary layer solver against the same Goldman case. The input Mach distribution was reconstructed from the published design parameters. The peak $H_i$ location and the separation verdict reproduce, while the peak magnitude overshoots (2.6 against Goldman's 2.2 at the same chord position).}
    \label{fig:goldman-validation-hi}
\end{figure}

\paragraph{Starting}

On startup before the turbine passage runs supersonically, a detached bow shock stands ahead of the leading edges. The flow behind this normal shock is subsonic and has lost a significant amount of total pressure. For the turbine to operate nominally with supersonic flow, the passsage must be able to "swallow" this shock, allowing the shock to move through the whole passage. \todo{make sure this explains at standstill}

To do so, the post shock flow with a reduced total pressure must have enough of a flow area pass the required mass flow similar to the choked flow condition. A similar condition is the unstart of a supersonic wind tunnel or a supersonic internal intake. However, Boxer notes that due to the large velocity gradient between the channel walls due to tue turning, the flow cannot be assumed one directional as in the previous examples. Goldman in 1968 provides the condition accounting for the "reduction of mass flow due to curvilinear flow" to calculate the maximum inlet design mach number to ensure starting as:

\begin{equation}
\frac{p_{02}}{p_{01}}(M_i) \;\geq\; \frac{Q}{1-C}
\label{eq:goldman_starting}
\end{equation}

\todo{what the fuck is this section}

Each term carries real content, so they are written out in full here because the off-design version of this inequality is what explains the coupled test in \cref{sec:results}. The left side is the total pressure recovery across the detached normal shock at the inlet design Mach number $M_i$, given by the Rankine and Hugoniot relation 

\begin{equation}
\frac{p_{02}}{p_{01}}(M_i) = \left[\frac{(\gamma+1)M_i^2}{(\gamma-1)M_i^2+2}\right]^{\frac{\gamma}{\gamma-1}}
\left[\frac{\gamma+1}{2\gamma M_i^2 - (\gamma-1)}\right]^{\frac{1}{\gamma-1}}
\label{eq:normal_shock_p0}
\end{equation}

which falls steeply as $M_i$ rises, so a higher inlet Mach number leaves less recovered pressure to push the swallowed shock through. The right side is the post shock area demand. $Q$ is the vortex flow parameter, the weight flow the curved passage would pass relative to a straight channel, written over the passage bounded by the lower and upper surface critical Mach numbers $M^*_l$ and $M^*_u$ (related to the true Mach numbers by \cref{eq:mach_to_critical_mach}) as

\begin{equation}
Q = \frac{M^*_l\, M^*_u}{M^*_u - M^*_l}
\int_{M^*_l}^{M^*_u}\!\left(\frac{\gamma+1}{2} - \frac{\gamma-1}{2}M^{*2}\right)^{\frac{1}{\gamma-1}}\frac{dM^*}{M^*}
\label{eq:vortex_parameter}
\end{equation}

and $C$ is the reduction in passable mass flow caused by the two dimensional velocity gradient across the passage that Boxer identified,

\begin{equation}
C = 1 - \sqrt{\frac{\gamma+1}{\gamma-1}}\left(\frac{\gamma+1}{2}\right)^{\frac{1}{\gamma-1}}
\frac{K^*_{max}\, M^*_u}{M^*_u - M^*_l}\, I_R, \qquad
I_R = \left(1 - K^{*2}_{max}\right)^{\frac{1}{\gamma-1}} - \left[1 - K^{*2}_{max}\!\left(\frac{M^*_u}{M^*_l}\right)^2\right]^{\frac{1}{\gamma-1}}
\label{eq:twod_reduction}
\end{equation}

where $K^*_{max}$ is the value that maximises the swallowed weight flow, found by solving the weight flow identity

\begin{equation}
\int_{M^*_l}^{M^*_u}\!\left[1 - \left(\frac{K^*_{max}}{M^*_l}\right)^2 M^{*2}\right]^{\frac{1}{\gamma-1}}\frac{dM^*}{M^*}
= \left(1 - K^{*2}_{max}\right)^{\frac{1}{\gamma-1}} - \left[1 - K^{*2}_{max}\!\left(\frac{M^*_u}{M^*_l}\right)^2\right]^{\frac{1}{\gamma-1}}
\label{eq:kstar_max}
\end{equation}

Solving \cref{eq:goldman_starting} as an equality for $M_i$ gives the maximum inlet design Mach number that will still start, which is the limit applied in the blade design of \cref{sec:design}. The two integrals \cref{eq:vortex_parameter,eq:kstar_max} have no closed form and are evaluated numerically. The physical reading is the one used later: the limit is set by the surface Mach pair $(M^*_l, M^*_u)$, so the same passage that is opened up for efficiency by spreading the surface Mach numbers is simultaneously pushed towards the starting limit, and at low blade speed the rising relative inlet Mach number crosses it. \todo{perhaps put in the starting graph, showing starting limit as a function of inlet and surface mach numbers?}

Historically, starting was a major problem for supersonic turbines. Before the theory was established, Colclough had to build a special jig in which every blade could be pivoted to vary the contraction until the blade started.

Colclough's cascade tests also produced something the inviscid theory cannot, an empirical derating of the starting limit that accounts for boundary layer blockage. His measured startable contraction is $k = A_{min}/A_2 = 1.029 - 0.163(M_2 - 1)$, valid for $1.0 < M_2 < 1.4$ \cite{colcloughDesignTurbineBlades1966}. This is applied on top of the Goldman limit in the design of \cref{sec:design}, and the design inlet relative Mach number of this turbine sits inside its validity range.

\paragraph{The efficiency trade}

The separation criterion limits the acceleration of the flow along the upper surface, and the deceleration of the flow along the lower surface. Then, the starting criterion provides an additional limit on the inlet design mach number based on these surface mach numbers. 

However, Goldman notes that for optimum efficiency, the maximum spread of surface number leads to maximum efficiencies. It comes from the obvious fact that a loss of velocity across the impulse blade leads to a loss of kinetic energy, so he posits the efficiency of an impulse blade to be:

\begin{equation}
    \eta = \left(\frac{w_{exit}}{w_{exit,ideal}}\right)^2
\end{equation}



This is nearly the same as Weiss' \cref{eq:weiss_rotor_coeff} loss. Weiss derived his empirical model only of 1D parameters such as inlet mach number and the angled the flow is turned. 

What this efficiency tries to capture is the boundary layer friction losses, shocks, and exit wake mixing losses. The effect of spreading the surface mach numbers becomes clear by looking at the vortex equation \cref{eq:vortex_condition_nondim}. The larger the mach number spread, the larger the difference in $R^*$ between the upper and lower surface, the larger physical passage width. The wider the channel compared to the blade chord, the lower fraction of the boundary layer thickness to passage width, the lower the wetted surface area per unit flow, and thus the lower the losses.

This is the challenge of supersonic turbine blade design, maximum efficiency is achieved by riding along the edge of separation and starting.

\begin{figure}[H]
    \centering
    \begin{subfigure}[b]{0.49\textwidth}
        \centering
        \includegraphics[width=\textwidth]{moc_spread_narrow.pdf}
        \caption{$M_l = 1.25$, $M_u = 1.35$}
        \label{fig:moc-spread-narrow}
    \end{subfigure}
    \hfill
    \begin{subfigure}[b]{0.49\textwidth}
        \centering
        \includegraphics[width=\textwidth]{moc_spread_wide.pdf}
        \caption{$M_l = 1.05$, $M_u = 1.5$}
        \label{fig:moc-spread-wide}
    \end{subfigure}
    \caption{The effect of surface Mach number spread on the blade passage, for the same inlet conditions and shared axis scales. A narrow spread produces a sliver of a passage where the boundary layers occupy a large fraction of the width, while a wide spread opens the passage up. This is why efficiency rises with Mach spread until separation or starting limits are hit.}
    \label{fig:moc-spread}
\end{figure}

\subsubsection{Finite Leading Edge}

In \cref{subsec:method_of_characteristics}, the upper and lower surface are mathematically designed to produce no shocks in transitioning the unifrom flow from the nozzle exit into the vortex flow. For no shocks to form, both surfaces must be perfectly parallel to the starting flow. This is problematic because this implies the point where the two surfaces intersect at the leading edge, the angle between the upper and lower surfaces is zero. This creates an impossible to manufacture zero thickness leading edge. Boxer suggests a few methods to solve this issue. He separates design into turbines with a subsonic or supersonic axial velocity component. 

For subsonic axial velocity turbines, either lower or upper surface can be designed with an incidence angle to the inlet flow. Angling the upper convex surface and keeping the lower concave surface parallel to the inlet flow, the compressive oblique shock will be generated along the region between the inlet and the nozzle exit, affecting the flow in the next blade passage. Kantrowitz shows that in steady state a compression shock of equal strength will be created to obtain the steady state condition \cite{kantrowitz}. To solve this, a small angle wedge can be placed shortly after the leading edge to generate a small expansion wave to cancel out the shock. 

Another method is angling the lower concave surface so that the shock stays within the blade passasge and does not affect the upstream flow. The oblique comrpession shock can then be designed to intersect the upper surface at an angle to fully cancel out the shock.

For supersonic axial velocity turbines, a wedge can be designed on the convex surface with a low enough angle from the design inlet flow so that the shock stays within the blade passage and can be cancelled out by a kink on the concave surface.

\todo{figures to explain this}

\begin{figure}[H]
    \centering
    
\includegraphics[width=0.6\textwidth]{placeholder.png}
    \caption{Finite leading edge treatments for supersonic blades: incidence angling of one surface, in passage shock cancellation, and the small wedge used to cancel the resulting compression shock.} % PLACEHOLDER IMAGE: hand sketch
    \label{fig:finite-le}
\end{figure}

\subsubsection{Nozzle Design}

The nozzle throat sets the stage mass flow. Once the pressure ratio exceeds the critical value the throat chokes, and the mass flow depends only on the throat area and the upstream stagnation state:

\begin{equation}
    \dot{m} = \frac{A_{th}\, p_{01}}{\sqrt{T_{01}}}
    \sqrt{\frac{\gamma}{R}} \left(\frac{2}{\gamma+1}\right)^{\frac{\gamma+1}{2(\gamma-1)}}
    \label{eq:choked_massflow}
\end{equation}

This equation does double duty in this thesis. In design it sizes the nozzle throats for the required mass flow, and in the analysis of \cref{sec:results} it is the only available estimate of the air mass flow, since the rig carries no air flowmeter. Its key property is that it is independent of shaft speed and of everything downstream of the throat, which makes it valid whether or not the rotor is operating as intended.

While the throat sets the mass flow, the ratio of the exit area to the throat area sets the exit Mach number, through the isentropic area relation

\begin{equation}
    \frac{A_3}{A_{th}} = \frac{1}{M_3}\left[\frac{2}{\gamma+1}\left(1 + \frac{\gamma-1}{2}M_3^2\right)\right]^{\frac{\gamma+1}{2(\gamma-1)}}
    \label{eq:area_mach}
\end{equation}

There is no closed form for $M_3$, so \cref{eq:area_mach} is inverted numerically for the supersonic root given the design area ratio. This relation is what links the nozzle geometry to the velocity triangle: the exit Mach number it returns must match the $M_3$ demanded by the inlet swirl of \cref{eq:triangle_inversion,eq:sizing_exit_state}, and the same relation is solved again when the design is re-evaluated for a different working gas, since the area ratio is fixed by the printed hardware but the Mach number and exit velocity are not.

Historically, the optimal nozzle design for a ideal inviscid gas was designed with the Method of Characteristics as well. Rao provides a method to achieve this \cite{raonozzledesign}. NASA special pubplication 8120 explains that a single canted parabola closely approximates the ideal nozzle contour and should be used for high performance nozzles. In this paper, the nozzle was a simple conical nozzle with a 7 degree half angle. It can be seen from \todo{add figure 7 from nasa sp8120} that the performance of a conical nozzle for low area ratios lead to a efficiency loss below 2\%.

\begin{figure}[H]
    \centering
    
\includegraphics[width=0.55\textwidth]{placeholder.png}
    \caption{Conical against contoured nozzle efficiency as a function of area ratio, reproduced from NASA SP-8120. At the low area ratio of this design the conical penalty is below 2\%.} % PLACEHOLDER IMAGE: SP-8120 Fig. 7 scan
    \label{fig:conical-nozzle-loss}
\end{figure}

\subsubsection{Turbopump Matching}

When the pump and turbine are coupled on one shaft with no speed control, the system settles where the torques balance:

\begin{equation}
    \tau_{\textrm{turbine}}(\omega) = \tau_{\textrm{pump}}(\omega) + \tau_{\textrm{friction}}(\omega)
    \label{eq:matching}
\end{equation}

The two sides have opposite slopes. The pump torque rises with speed, approximately with $\omega^2$ from the affinity laws. The turbine torque falls with speed, because for a fixed supply the jet velocity is fixed, so increasing blade speed reduces the velocity difference that generates the blade force. The intersection of the two curves is the match point. Because the turbine curve falls and the load curve rises, a small disturbance in speed produces a net torque that pushes the shaft back towards the intersection, so the match point is self regulating and no active speed control is needed. This is the behaviour the coupled test in \cref{sec:results} is designed to demonstrate.
